\section{Markowitz experiment: details}\label{app:mk}
% =============================================================================================
% Owner: markowitz writer. Completes sec:markowitz (E4, experiments/e4_markowitz.py; as run:
% 30 universes, B = 20, 2 test seeds, 60 placebo replications; results/e4/meta*.json).
% Generated tables: e4_trtr, e4_utility_agg, e4_tests (+ note), e4_placebo (+ note),
% e4_link (+ note), e4_link_pooled; figure e4_placebo_size. Numbers not in a generated table
% (fit diagnostics, counts over universes, the Kantorovich check, the sign count of the link
% table) were computed from results/e4/{e4_diagnostics,e4_utility,e4_trtr,e4_link_points,
% e4_link}.csv with short scripts. The residual-correlation null sqrt(1/Ttr + 1/(N-1)^2) = 0.0445
% (Ttr = 1258, N = 30; the code's resid_corr_rms_null, FactorGARCHt.fit_info_) and the ratio of
% resid_corr_rms to it (2.52 to 3.88 over universes) were computed from e4_diagnostics.csv.
% Seed-to-seed reliability: tab:manual-e4-reliability in sec:mk-link (paper/scripts/
% e4_reliability.py from results/e4/e4_tests.csv; values re-checked here from the same CSV).
% The 18 of 192 link combinations (4 positive; smallest generator-level exact p = 0.083) were
% recounted from e4_link.csv; the chance count 8.0 = 0.05 x 144 + (2/120) x 48. Attenuation
% bounds 0.50 and 0.78: sqrt(2r/(1+r)) with the mean within-generator r of
% tab:manual-e4-reliability (0.14 for T, 0.44 for MMD^2_u); ten seeds (DESIGN): 10r/(1+9r) = 0.62
% for T. Critic training windows per sample: 123 - n_eval = 60 to 63 (e4_tests.csv).
% =============================================================================================

This appendix completes \cref{sec:markowitz}. It reports fit diagnostics of the five generators,
the TRTR benchmark and the three aggregations of the utility losses, the full test table with
the placebo that calibrates it, the full link tables and what the seed-to-seed reliability of the
test quantities implies for the link. Numbers that do not appear in a table
were computed from the E4 output files. All results use the settings as run: 30 universes of 30
assets, $B=20$ synthetic replicas, two synthetic seeds in the tests part and 60 placebo
replications per generator.

% ---------------------------------------------------------------------------------------------
\subsection{Generators and diagnostics}\label{app:mk-generators}

The fitted degrees of freedom of MV Student-$t$ have median 5.08 across universes (range 4.57 to
5.67). The MV bootstrap has a median mean block length of 52.9 days (range 50.9 to 54.4). For
CCC-GARCH-$t$ the mean off-diagonal copula correlation is 0.356 on average over universes. Part
of the attenuation of its return correlations thus comes from the copula estimate (0.356 against
0.419 in the data) and part from the dynamics (0.310 in simulated samples;
\cref{tab:e4-correlation}). Its margins have a median
persistence of 0.968 (mean over universes of the per-universe median), and 14 of the 900 margins
reached the persistence cap of 0.995.

For factor GARCH-$t$ the single-index regressions $r_{it}=a_i+b_im_t+e_{it}$ have a mean $R^2$ of
0.447 on average over universes (range 0.392 to 0.492). The factor $m_t$ is the equal-weight mean
of the log-returns of the same $N=30$ assets, so the OLS residuals add up to zero every day,
$\sum_ie_{it}=0$. Hence $\sum_{i\ne j}\Cov(e_i,e_j)=-\sum_i\Var(e_i)$ in sample, which holds the
average pairwise residual correlation close to $-1/(N-1)$ whatever dependence is left in the
residuals, so that average is uninformative. The root mean square of the pairwise residual
correlations is informative. It is 0.140 on average over universes (range 0.112 to 0.173), against
$\sqrt{1/\Ttr+1/(N-1)^2}=0.045$ for residuals that are independent across assets, as the model
assumes; this null is sampling noise of order $\Ttr^{-1/2}$ plus the adding-up term. The excess
over the null (the ratio is 2.5 to 3.9 across universes) measures the dependence beyond the market
factor, such as industry effects, that the model leaves out by setting the residual correlations
to zero.

The adding-up identity also shapes the simulated covariance. With variance targeting, the
unconditional covariance of the model, $bb'\Var(m)+\diag\{\Var(e_1),\dots,\Var(e_N)\}$ with
$b=(b_1,\dots,b_N)'$, equals the \TRAIN{} sample covariance (divisor $\Ttr$) except for the
off-diagonal residual covariances, which it sets to zero. It therefore drops two things: the
positive residual covariances of industry effects and the negative ones that the adding-up
identity creates. In the data the equal-weight mean return is the factor itself. In simulated
samples the residuals are independent and do not add up to zero, so the equal-weight mean also
carries the idiosyncratic variance $N^{-2}\sum_i\Var(e_i)$. Both parts change the low-variance
directions on which GMV portfolios load. E4 does not separate their contributions to the GMV loss
of factor GARCH-$t$, and the variant with an external factor, which would remove the adding-up
identity, was not run (\cref{sec:mk-utility}). Twelve of its 930 GARCH components reached the
cap. No portfolio optimisation failed in any replica.

The Kantorovich bound of \cref{prop:gmv-sensitivity}(b) was checked replica by replica, with
$S_{\mathrm{te}}$ the \TEST{} sample covariance of daily log-returns, the matrix that also enters
$\kappa$ in \cref{tab:e4-correlation} (the variance ratios of \cref{tab:e4-utility} use the daily
simple returns of the portfolios instead): for every replica-wise GMV portfolio (sample moments),
the variance under $S_{\mathrm{te}}$ relative to that of the ex-post GMV portfolio
$\wgmv(S_{\mathrm{te}})$ was below $(1+\kappa_G)^2/(4\kappa_G)$. The bound is loose. The median
bound of the real \TRAIN{} covariance is 5.177 (\cref{tab:e4-correlation}), while the TRTR GMV
portfolio has 1.58 times the variance of the ex-post GMV portfolio on average over universes.

% ---------------------------------------------------------------------------------------------
\subsection{TRTR benchmark and aggregation of the losses}\label{app:mk-aggregation}

\begin{table}[tbp]
  \centering
  \caption{Out-of-sample performance of the portfolios trained on real \TRAIN{} (TRTR) on the 615
  \TEST{} days (E4). Means over the 30 universes. Vol.: annualised standard deviation of daily
  simple portfolio returns; Sharpe: annualised Sharpe ratio with a zero risk-free rate; CE:
  certainty-equivalent return with the $\gamma$ of the rule ($\gamma=5$ for GMV and $1/N$),
  percent per year. Sample: sample covariance; LW: Ledoit--Wolf. Point estimates.}
  \label{tab:e4-trtr}
  \small
  \fittable{e4_trtr}
\end{table}

\Cref{tab:e4-trtr} gives the benchmark against which every loss is measured. Ledoit--Wolf
shrinkage lowers the volatility of the GMV portfolio slightly (13.89\% against 14.12\% per year).
The long-only GMV portfolio has a lower volatility than the unconstrained one (13.33\%), the
familiar effect of the sign constraint as shrinkage \citep{jagannathan2003risk}. The $1/N$
portfolio has the highest Sharpe ratio, 0.83, and a higher certainty equivalent at $\gamma=5$ than
the mean-variance portfolio with $\gamma=5$ (6.97\% against 4.64\% per year), as in
\citet{demiguel2009optimal}. The CE values of rules with different $\gamma$ are not comparable
with each other. A positive risk-free rate would lower the Sharpe ratios, and could change their
ranking across portfolios of different volatility, but not the CE comparison of $1/N$ with the
mean-variance portfolio, because every portfolio is fully invested (\cref{sec:mk-design}).

\begin{table}[tbp]
  \centering
  \caption{Utility losses under three aggregations over the $B=20$ replicas (E4; sample moments;
  means over 30 universes). Replica: mean loss of the 20 portfolios built on single replicas (as
  in \cref{tab:e4-utility}); Pooled: one portfolio from the moments of the 20 stacked replicas;
  Avg.\ $w$: the mean of the 20 replica-wise weight vectors. Var.\ ratio: \TEST{} variance of the
  TSTR portfolio over the TRTR one; CE loss: TRTR minus TSTR certainty equivalent, percent per
  year. Point estimates.}
  \label{tab:e4-utility-agg}
  \small
  \fittable{e4_utility_agg}
\end{table}

\Cref{tab:e4-utility-agg} separates misspecification from estimation noise. Pooling the replicas
lowers the variance ratio for every generator and rule except the unconstrained GMV portfolio of
factor GARCH-$t$ (1.436 replica-wise, 1.466 pooled). For the mean-variance rules it lowers the CE
loss in 14 of the 15 cases, the exception being CCC-GARCH-$t$ with $\gamma=1$
(\cref{cor:gmv-drivers}(iii)). The averaged weights satisfy
\cref{prop:no-free-lunch}: in each of the 1500 combinations of universe, generator, estimator and
rule ($1/N$ excluded), both the variance ratio and the CE loss of the averaged weights are at
most the mean of the replica-wise ones. For the long-only mean-variance rules with $\gamma=1$ the
averaged weights have variance ratios between 0.606 and 0.706: averaging concentrated long-only
portfolios over replicas diversifies them, at a CE cost that remains positive for every generator.

% ---------------------------------------------------------------------------------------------
\subsection{Tests and placebo}\label{app:mk-placebo}

\begin{table}[p]
  \centering
  % Landscape: the table is rotated with \rotatebox (graphicx), no extra package.
  \rotatebox{90}{\begin{minipage}{0.95\textheight}
    \centering
    \caption{Multivariate tests of E4 with Wilson 95\% intervals: full version of
    \cref{tab:e4-tests}. $\alpha=0.05$; 30 universes; two synthetic seeds.}
    \label{tab:e4-tests-full}
    \footnotesize
    \fittable{e4_tests}
    \tablenote{\inputnote{e4_tests_note}}
  \end{minipage}}
\end{table}

\begin{table}[p]
  \centering
  % Landscape: the table is rotated with \rotatebox (graphicx), no extra package.
  \rotatebox{90}{\begin{minipage}{0.95\textheight}
    \centering
    \caption{Placebo size of the E4 tests under $H_0$: both samples are independent \TEST{}-length
    draws from the generator fitted on one universe, tested exactly as in \cref{tab:e4-tests} (60
    replications per generator; rejection rates with Wilson 95\% intervals at three levels).}
    \label{tab:e4-placebo}
    \footnotesize
    \fittable{e4_placebo}
    \tablenote{\inputnote{e4_placebo_note}}
  \end{minipage}}
\end{table}

\begin{figure}[tbp]
  \centering
  \inputfigure{e4_placebo_size}
  \caption{Placebo rejection rates of the E4 tests against the nominal level $\alpha\in[0,0.2]$:
  empirical distribution functions of the HAC p-values of the \WC{} test (left) and of the
  wild-bootstrap p-values of MMD (right), 60 placebo replications per generator. Dashed: the
  diagonal (exact size); dotted: $\alpha=0.05$.}
  \label{fig:e4-placebo-size}
\end{figure}

\paragraph{Placebo design.}
For each generator, placebo replication $r=0,\dots,59$ uses the generator fitted on universe
$r\bmod30$, so each universe contributes two replications. Two independent \TEST{}-length samples
$X^*$ and $Y^*$ are simulated from the end of \TRAIN{}, divided by the real \TRAIN{} standard
deviations, cut into $(30,5)$ windows with \TEST{} date ranks and tested with exactly the call of
the tests part; the MMD bandwidth is the median heuristic of the $X^*$ windows, as it is that of
the real \TEST{} windows in the tests part. Each replication records the p-values and a score:
$\zhac$ for the \WC{} test and $-\pwild$ for MMD (larger is more evidence against $H_0$).

\paragraph{Results.}
\Cref{tab:e4-tests-full} repeats \cref{tab:e4-tests} with intervals and the mean statistics.
\Cref{tab:e4-placebo} and
\cref{fig:e4-placebo-size} give the placebo. For the two \iid{} generators the Driscoll--Kraay
p-value of the \WC{} test has its nominal size (0.050 and 0.067 at 5\%), the studentised
statistic $\zhac$ has mean $-0.13$ and $0.05$ and standard deviation $0.98$ and $1.03$, and the KS
test does not reject uniformity of the p-values. For the generators with persistent volatility
the test is anti-conservative. For CCC-GARCH-$t$, factor GARCH-$t$ and the MV bootstrap, in this
order, the size at 5\% is 0.267 [0.171, 0.390], 0.267 [0.171, 0.390] and 0.133 [0.069, 0.242], and
$\zhac$ has mean 0.87, 0.35 and 0.22 and standard deviation 1.34, 1.47 and 1.37. The size of the
wild bootstrap of MMD also exceeds the nominal level for these three generators (0.117, 0.100 and
0.133), and the permutation p-values, which ignore the dependence between windows, are oversized
for all three under both tests.

The ordering, nominal for \iid{} generators, intermediate for the block bootstrap (mean blocks of
about 53 days) and worst for the GARCH generators, points to volatility persistence. Two
mechanisms are consistent with the moments of $\zhac$. Its positive mean indicates that $T$ is
biased upwards under $H_0$: within one sample, training and evaluation blocks are adjacent pieces
of the same persistent path, with no gap between them, so a critic that learns the volatility
regime of one sample in its training blocks still finds it in the adjacent evaluation blocks.
\Cref{lem:decoupling} controls this dependence only with a gap, which the time-block split does
not have (\cref{sec:method-inference}). Its standard deviation above one indicates that the
Driscoll--Kraay variance is too small with 60 to 63 evaluation dates (one per window) per
sample (\cref{rem:small-T}). The 123 windows of a 615-day sample are few; the E3 panel has a different
geometry (100 tickers per date) and its size is measured separately in E1.

The size-adjusted shares of \cref{tab:e4-tests} use the placebo as a Monte Carlo null. The
calibrated p-value $p_{\mathrm{cal}}=(1+n_s)/61$, where $n_s$ is the number of the 60 placebo
scores of the same generator and test at least as large as the score of seed $s$, is a Monte
Carlo p-value for the hypothesis that the real \TEST{} sample is a draw from the fitted
generator. It treats the universe as drawn at random (the placebo pools two replications from
each universe), so it is approximate for a given universe. It is also discrete. With two seeds,
$\pmed=\min\{1,p_{\mathrm{cal},1}+p_{\mathrm{cal},2}\}$ is at least $2/61\approx0.033$, and it is at
most 0.05 only if $n_1+n_2\le1$: both seeds must score above all, or all but one, of the placebo
replications. This explains why the size-adjusted MMD shares with $\pmed$ are zero for the three
dynamic generators while the first-seed shares are not. The reference row (last 615 days of real
\TRAIN{}) has no placebo, and its nominal rejection share cannot be corrected for the oversize
documented here.

% ---------------------------------------------------------------------------------------------
\subsection{Link between tests and utility}\label{app:mk-link}

\begin{table}[p]
  \centering
  % Landscape: the table is rotated with \rotatebox (graphicx), no extra package.
  \rotatebox{90}{\begin{minipage}{0.95\textheight}
    \centering
    \caption{Spearman correlations between test quantities and utility losses (E4; replica-wise
    portfolios, sample moments; full version of panel A of \cref{tab:e4-link}). Intervals: 2000
    universe-bootstrap draws; $p_u$, $p_g$: 2000 permutations.}
    \label{tab:e4-link-full}
    \footnotesize
    \fittable{e4_link}
    \tablenote{\inputnote{e4_link_note}}
  \end{minipage}}
\end{table}

\begin{table}[p]
  \centering
  % Landscape: the table is rotated with \rotatebox (graphicx), no extra package.
  \rotatebox{90}{\begin{minipage}{0.95\textheight}
    \centering
    \caption{As \cref{tab:e4-link-full}, for portfolios built from the moments of the 20 pooled
    replicas, which isolate the misspecification of the generator (full version of panel B of
    \cref{tab:e4-link}). Intervals and p-values as in \cref{tab:e4-link-full}.}
    \label{tab:e4-link-pooled}
    \footnotesize
    \fittable{e4_link_pooled}
  \end{minipage}}
\end{table}

\paragraph{Link design.}
Each of the 150 (universe, generator) points carries the seed means of $T$, $\zhac$, $\TMMD$ and
the wild-standardised MMD statistic (the MMD V-statistic centred and scaled by the mean and
standard deviation of its wild-bootstrap replicates), the correlation error
$\abs{\Delta\,\text{avg.\ corr.}}$, and the losses of every rule, estimator and aggregation. The
studentised statistics replace the p-values, which lose resolution for strong rejections. The
test quantities come from one \TEST{}-length sample per seed started at the end of \TRAIN{} and
the losses from $B$ \TRAIN{}-length samples from a stationary start, so the link concerns the
unconditional moments of each generator (\cref{sec:mk-design}). Four
views are reported because there are only five generators. The \emph{pooled} Spearman correlation
uses all 150 points. The \emph{within-universe} correlation is computed across the five
generators of each universe and averaged over universes; its permutation p-value permutes the
generator labels of the losses within each universe, so it is conditional on these generators.
Both views mostly rank the five generators. The \emph{within-generator} correlation is computed
across the 30 universes of each generator and averaged over generators; its permutation p-value
permutes universes within each generator, and it is insensitive to generator-level offsets of a
statistic, such as the placebo bias of the critic. The \emph{generator-level} correlation is that of
the five generator means, with an exact p-value over the 120 orderings; it is the only view whose
sampling unit is the generator. Intervals resample universes with replacement, keeping all
generators of a drawn universe together.

\paragraph{Results.}
\Cref{tab:e4-link-full,tab:e4-link-pooled} give all six losses. Besides the GMV variance ratio,
the absolute correlation error is associated with the long-only GMV loss when the replicas are
pooled ($-0.37$ [$-0.55$, $-0.19$] pooled, $-0.21$ within generator, $p=0.014$), more weakly than
for the unconstrained portfolio, plausibly because the sign constraint limits the long-short
positions through which correlations act. As for the unconstrained portfolio, the negative sign
reflects an evaluation period in which correlations fell (\cref{sec:mk-link}). Neither the test
quantities nor the correlation error has a consistent association with any mean-variance loss.
Counting over both tables, 18 of the 192 combinations of the four test quantities, six losses,
two aggregations and four views have a pooled interval that excludes zero or a permutation
p-value below 0.05; the four positive ones are $T$ against the pooled-moment GMV and long-only GMV
losses, in the pooled and within-universe views, and the 14 others are negative and all involve
a studentised statistic ($\abs{\rho}$ between 0.15 and 0.31). The count is descriptive. If all
associations were null, the comparisons independent and the intervals exact, about eight would
be expected by chance, slightly fewer than 5\% of 192 (\cref{sec:mk-link}): 0.05 for each of the
144 comparisons of the first three views and $2/120$ for each of the 48 generator-level ones,
none of which reaches 0.05 here. The comparisons share the 150 points and the intervals are
optimistic, so the count varies far more under the null than a binomial count would. The one
pattern that recurs, the positive association of $T$ with the pooled-moment GMV losses, is
between generators: within generators it vanishes.

The generator-level view has $n=5$. Its exact two-sided p-value cannot fall below
$2/120\approx0.017$, which requires a perfect rank agreement or reversal; the largest
generator-level correlations in the tables, $0.90$ for $T$ and $-0.90$ for the correlation error,
both against the Sharpe loss of the mean-variance portfolio with $\gamma=5$, have $p=0.083$. Any
statement about which kind of generator is more useful therefore rests on five generators and
does not extend beyond them. The within-generator view is the one that a single generator's user
faces: across universes, it shows no association between any test quantity and the
unconstrained or long-only GMV loss, under either aggregation (all permutation p-values above
0.05; \cref{tab:e4-link-full,tab:e4-link-pooled}).

\paragraph{Reliability of the test quantities.}
How much this absence says depends on how reproducible the test quantities are across synthetic
seeds (\cref{tab:manual-e4-reliability}). Within generators, the mean seed-to-seed Spearman
correlation across universes is 0.14 for $T$, 0.26 for $\zhac$ and 0.44 for $\TMMD$. Under a
classical measurement model, the seed mean of two runs, which enters the link, has reliability
$2r/(1+r)$ for a seed-to-seed correlation $r$, and its correlation with any fixed property of a
universe, such as a loss, is at most the square root of that reliability: about 0.50 for $T$ and
0.78 for $\TMMD$ at these mean values of $r$. A moderate association between the expected critic
statistic and the GMV loss would thus be at least halved and hard to detect with 30 universes
per generator, so the within-generator correlations of $T$ with the GMV losses (0.01
replica-wise, $-0.01$ with pooled moments) are weak evidence of no association. The same absence
is more informative for MMD, whose statistic is more reproducible, most for the two \iid{}
generators. The reference row, whose data are the same under both seeds, still has a
seed-to-seed correlation of 0.28 for $T$, so the test itself, through the random split and the
training of the critic on 60 to 63 windows per sample, contributes much of this noise. The ten
seeds of the design would raise the reliability of the seed mean of $T$ to about 0.62 under the
same model; two were run. No architecture or training ablation of the multivariate critic was
run either. The evidence that the tests and utility order the generators differently therefore
rests on the comparison of generators, with the caveat of $n=5$ above, and on MMD
(\cref{sec:mk-link}).
